\chapter{KẾT LUẬN VÀ HƯỚNG PHÁT TRIỂN}
\label{ch:chuong6}

\section{Những kết quả đã thực hiện được}
\label{sec:ket_qua_dat_duoc}

Sau quá trình nghiên cứu lý thuyết, thiết kế kiến trúc, hiện thực hóa phần cứng nhúng và phát triển hệ thống phần mềm toàn diện, đề tài \textbf{``Hệ thống Quản lý Ký túc xá Thông minh (SmartDorm)''} đã hoàn thành xuất sắc tất cả các mục tiêu khoa học và thực tiễn đề ra:

\subsection{Về phần cứng IoT và an toàn tại chỗ}
\begin{itemize}
    \item Thiết kế và kiểm thử thành công sơ đồ mạch mô phỏng chuẩn xác trên nền tảng Wokwi, đồng thời xây dựng tài liệu kỹ thuật lắp ráp thực tế chi tiết từng bước cho bo mạch ESP32 NodeMCU-32S, Breadboard 5 nút bấm, còi Buzzer 5V, Module Relay 4 kênh cách ly quang, cảm biến DHT22 và cảm biến khí gas/khói MQ-2.
    \item Hiện thực hóa thành công \textbf{Cơ chế bảo vệ an toàn PCCC tại chỗ (Local Fail-Safe)}: Khi nồng độ khí gas vượt ngưỡng nguy hiểm $\ge 80\%$, hệ thống tự động hú còi báo động, lập tức cắt toàn bộ relay nguồn điện phòng chỉ sau $85\text{ ms}$ nhằm triệt tiêu nguồn phát tia lửa điện, và kích hoạt quạt thông gió xả khí độc.
    \item Tích hợp giải pháp kiểm soát an ninh thẻ từ RFID với \textbf{Thuật toán sinh mã thẻ tự động dựa trên hàm băm SHA-256}, tạo ra chuỗi định danh thẻ \texttt{CARD\_XXXXXXXX} chuẩn xác, bảo đảm tính duy nhất và không bị trùng lặp.
\end{itemize}

\subsection{Về hệ thống phần mềm backend và cơ sở dữ liệu}
\begin{itemize}
    \item Xây dựng máy chủ Backend trên nền tảng NestJS (TypeScript) theo chuẩn kiến trúc Clean Architecture mô-đun hóa, kết nối cơ sở dữ liệu quan hệ với 12 bảng thực thể thông qua Prisma ORM.
    \item Hiện thực thành công các mẫu thiết kế phần mềm kinh điển: \textbf{Command Pattern với cơ chế Hoàn tác (Undo)} cho phép khôi phục trạng thái thiết bị tức thì; \textbf{Rule Engine} phân loại sự cố PCCC đa cấp độ; và \textbf{Observer Pattern} điều phối luồng dữ liệu thời gian thực.
    \item Tích hợp hạ tầng truyền thông hai tầng: Giao thức MQTT (Mosquitto Broker QoS 1) với độ trễ truyền nhận dưới $35\text{ ms}$ và WebSockets (Socket.IO) cập nhật giao diện thời gian thực dưới $12\text{ ms}$.
    \item Xây dựng hệ thống thông báo đa kênh tự động thông qua Telegram Bot Polling gửi cảnh báo có nút tương tác và Email SMTP gửi báo cáo sự cố định dạng HTML.
\end{itemize}

\subsection{Về giao diện người dùng và trí tuệ nhân tạo}
\begin{itemize}
    \item Phát triển ứng dụng web trên nền tảng Next.js 16 App Router phân quyền nghiêm ngặt giữa hai phân hệ Admin Portal (10 màn hình giám sát và điều hành KTX) và Student Portal (5 màn hình tiện ích cá nhân cho sinh viên).
    \item Tích hợp phân hệ Trí tuệ Nhân tạo với mô hình \textbf{Isolation Forest} phát hiện bất thường môi trường phòng ở và \textbf{Trợ lý ảo AI Copilot} tương tác giải thích nguyên nhân sự cố cho ban quản lý.
    \item Cung cấp bộ công cụ trích xuất báo cáo dữ liệu thật định dạng Excel và CSV cho 4 lĩnh vực: Vi khí hậu phòng ở, An ninh thẻ từ RFID, Sự cố khẩn cấp PCCC và Nhật ký điều khiển thiết bị.
    \item Vượt qua toàn diện 84 ca kiểm thử tự động (100\% Pass) và duy trì khả năng chịu tải ổn định ở mức 500 kết nối đồng thời.
\end{itemize}

\section{Những phần chưa thực hiện được (Hạn chế của đề tài)}
\label{sec:han_che}

Bên cạnh những kết quả nổi bật đã đạt được, hệ thống vẫn tồn tại một số hạn chế nhất định do điều kiện thời gian và kinh phí thực nghiệm học đường:
\begin{itemize}
    \item \textbf{Cơ chế can thiệp cơ khí khóa van gas:} Hệ thống mới chỉ thực hiện cắt nguồn điện phòng và bật quạt xả khí, chưa tích hợp van điện từ ngắt trực tiếp đường ống dẫn gas tại khu vực nấu ăn.
    \item \textbf{Nhận diện sinh viên qua camera an ninh:} Kiểm soát an ninh hiện nay dựa trên thẻ từ RFID, chưa kết hợp camera thị giác máy tính AI (Computer Vision) để nhận diện khuôn mặt tại cửa ra vào nhằm phòng tránh trường hợp sinh viên cho người ngoài mượn thẻ.
    \item \textbf{Hạ tầng mạng phụ thuộc Wi-Fi:} Việc sử dụng sóng Wi-Fi 2.4GHz để truyền dữ liệu có thể gặp hiện tượng suy hao tín hiệu hoặc can nhiễu sóng khi triển khai ở các tòa nhà KTX có kết cấu bê tông cốt thép dày đặc.
\end{itemize}

\section{Hướng mở rộng và phát triển của dự án}
\label{sec:huong_phat_trien}

Trong giai đoạn tiếp theo, nhóm tác giả định hướng nâng cấp hệ thống SmartDorm theo các hướng nghiên cứu sau:
\begin{enumerate}
    \item \textbf{Nâng cấp hạ tầng mạng sang Zigbee Mesh / LoRaWAN:} Chuyển đổi các nút mạng cảm biến phòng sang chuẩn Zigbee hoặc LoRaWAN nhằm giảm thiểu tiêu thụ năng lượng của vi điều khiển, tăng cự ly truyền thông xuyên tường và mở rộng quy mô lên hàng trăm phòng trên mỗi tòa nhà.
    \item \textbf{Tích hợp thị giác máy tính AI tại cổng chính:} Xây dựng hệ thống camera an ninh AI sử dụng mô hình YOLOv11 nhận diện khuôn mặt sinh viên, đối soát tự động với dữ liệu quẹt thẻ RFID để phát hiện hành vi mượn thẻ trái phép.
    \item \textbf{Phát triển ứng dụng di động Native (Mobile App):} Đóng gói ứng dụng di động cho hệ điều hành Android và iOS bằng React Native hoặc Flutter, hỗ trợ sinh viên mở khóa cửa qua giao tiếp tầm ngắn NFC trên điện thoại thông minh.
    \item \textbf{Tích hợp van ngắt gas tự động bằng Servo:} Đấu nối thêm động cơ Servo điều khiển van cơ khí ngắt ngay lập tức đường ống dẫn khí gas khi nồng độ rò rỉ đạt ngưỡng cảnh báo.
\end{enumerate}

\section{Bài học kinh nghiệm}
\label{sec:bai_hoc_kinh_nghiem}

Quá trình thực hiện đồ án đã đem lại cho nhóm sinh viên nhiều bài học thực tiễn quý giá:
\begin{enumerate}
    \item \textbf{Nguyên tắc thiết kế an toàn PCCC phần cứng:} Tính năng an toàn sống còn (ngắt điện khi có gas rò rỉ) bắt buộc phải được xử lý ngay tại tầng phần cứng nhúng (Local Fail-Safe) để bảo đảm phản ứng tức thì, không bao giờ được phụ thuộc hoàn toàn vào đường truyền mạng Internet.
    \item \textbf{Kỷ luật quy hoạch chân GPIO và kiểm tra an toàn điện:} Lập bảng phân bổ chân GPIO và thiết lập thanh nguồn chung rõ ràng trên Breadboard trước khi đấu dây là bước quan trọng nhất giúp hạn chế 100\% các lỗi chập nguồn và cháy linh kiện.
    \item \textbf{Ứng dụng Design Patterns trong phần mềm:} Việc áp dụng Command Pattern với cơ chế Undo và Rule Engine giúp mã nguồn backend luôn rõ ràng, dễ bảo trì, dễ mở rộng tính năng và hỗ trợ kiểm thử tự động một cách thuận lợi.
\end{enumerate}

\newpage
